% intro.tex : Introduction (English manuscript, body text only, no preamble)
% 작성 2026-10-04. 명세: paper/manuscript/MANUSCRIPT_SPEC.md 3절(I1-I4, 700단어 목표).
% 인용 키: paper/manuscript/en/refs_intro.bib (refs_discussion.bib 와 같은 FirstAuthorYear 규칙).
%
% 수치 출처(원고 조립 때 shared/numbers.tex 매크로로 바꾼다. 자릿수는 지침 4.5):
%   30 / 13,606 / 343 ............ paper/claims/D_data_and_design/README.md 2절 E14, E10
%   1.31 / 2.51 / 11.03 .......... 같은 README E38 (fig1_z_summary.csv E_mean, 셀 E 기하 평균)
%   1.6 [1.0, 2.0], P = 0.001 .... C1_label0_safety/README.md 2.3 첫 행(셀 가중 -1.63 [-1.98, -1.04])
%   2.5 [1.9, 3.0], P = 0.001 .... C2_bias_diagnosis/README.md 2.1 A6(-2.45 [-3.04, -1.88]), 지역 행 1/4 은 C3 README E36a
%   거의 전부(약 99%) ............. C8_gain_source/README.md 1.3, B4(알래스카 지역 내, 결과 열람 뒤 설계)
%   0.29 degC(2007-2016) ......... Biskaborn 2019 초록(로컬 PDF 대조)
%   68 지점, 무작위 70/30 1회, 시험 R2 0.24 대 0.54 ... Gautam 2025 본문(로컬 PDF 대조)
%   35 years ..................... Streletskiy 2025 자료 제목(DataCite 대조)
% 열린 결정: 지역 영문 표기(W Russia, E Russia, Lena Delta)는 MANUSCRIPT_SPEC 13.1 의 7번 결정 대기.
% 새 실험(XA-XJ) 자리표시: 명세 11절 등록부에 서론 자리가 없어 넣지 않았다.

Active-layer thickness (ALT) is the maximum seasonal thaw depth of the ground above permafrost, and its change affects infrastructure built on permafrost\cite{Karjalainen2019,Ran2022b}. Permafrost temperature increased by 0.29\,$^{\circ}$C on global average between 2007 and 2016\cite{Biskaborn2019}. ALT is measured at a limited number of sites, many of them in the 35-year Circumpolar Active Layer Monitoring (CALM) compilation\cite{Streletskiy2025}. In the five regions evaluated here, the number of measured ALT values (labels) ranges from 30 in E Russia to 13,606 in Alaska, and the Alaskan labels lie at 343 locations of 1\,km (Table~1). The Stefan solution\cite{Stefan1891} predicts ALT as a coefficient $E$ times the square root of the thawing degree-days (TDD), with $E$ calibrated by local thaw-depth measurements\cite{Nelson1997}. The geometric mean of $E$ over labelled cells is 1.31\,cm\,($^{\circ}$C\,d)$^{-1/2}$ in the Lena Delta, 2.51 in W Russia and 11.03 on the Tibetan Plateau (Fig.~1b). However, because measurements in a new permafrost region (the target region) are few, the error of ALT maps there depends on models fitted to other regions (the source data).

Outputs of the Stefan or Kudryavtsev models have served as machine-learning (ML) input features\cite{Pilyugina2025,WangG2025}, $E$ has been predicted from covariates\cite{Ran2022b,ZhangC2024}, labels generated by physical rules have trained a zero-curtain detector\cite{Gay2026}, and process-model pre-training and residual modelling are used in other environmental fields\cite{Read2019,Jia2021,Willard2022}. Each component has precedents, and the studies differ in how they were evaluated. Gautam et al. compared a random forest with a Stefan model at 68 CALM sites in Alaska using one random 70/30 split and reported test coefficients of determination of 0.24 and 0.54, respectively\cite{Gautam2025}. Large-scale ALT maps have been validated with distance blocks\cite{Aalto2018,Karjalainen2019,Ran2022a}. For clustered spatial data, random cross-validation can overstate map accuracy\cite{Roberts2017,Ploton2020,MeyerPebesma2022}, whereas Wadoux et al. argued that unbiased map accuracy requires probability sampling with design-based inference rather than spatial cross-validation\cite{Wadoux2021}. The ALT evaluations above each use one label condition, whereas adjacent fields have evaluated models in held-out regions or with few local observations, in some cases against calibrated process models, for lake temperature\cite{Read2019,Jia2021,Willard2021}, ungauged catchments\cite{Pool2019,Feng2023}, subsurface temperature\cite{OMalley2026} and spatial risk prediction\cite{Portes2026}. To our knowledge, no study of permafrost ALT or mean annual ground temperature has reported the error in regions excluded from the training data as a function of the number of target labels, relative to a Stefan model recalibrated with those labels.

A new region needs a rule for combining the Stefan model and ML as target labels accumulate, and a procedure for adding labels. We address three questions. (1) When no target labels are available, what does physics information add to ML, and where does the gain come from? (2) As target labels increase from ten to thousands, how much does ML add to a Stefan model recalibrated with the same labels, and at which spatial scale does the gain arise? (3) As labels accumulate, how should the method be selected, and by which procedure should new labels be placed? The answers define the steps of a workflow.

Each target region was held out from the source data with a 100\,km buffer, labels were drawn from half of its 0.5$^{\circ}$ blocks, and errors were scored on the other half, for 0, 3, 10, 40, 160, 320, 1000 and all available labels (Fig.~1c,d). Every method was compared with the Stefan model before and after recalibration on the same labels, as in lake-temperature studies\cite{Read2019,Jia2021}, with the recalibrated coefficient shrunk towards the source value\cite{Steyerberg2004}. Within label-rich regions, block-holdout tests designed after earlier results were viewed measured the gain of ML over the recalibrated model. Contrasts were judged from block-bootstrap 95\% confidence intervals (CIs) with a $\pm$0.5\,cm equivalence margin, and the transfer analysis was internally pre-registered. We find that, without target labels, Stefan pseudo-labels (Stefan predictions used as training labels) reduced the mean root-mean-square error (RMSE) of four main regions (Lena Delta, Canada, W Russia and E Russia) relative to shuffled pseudo-labels by 1.6\,cm (95\% CI 1.0 to 2.0; Holm-adjusted $P = 0.001$), and that recalibrating the Stefan coefficient with ten labels reduced it by 2.5\,cm relative to the source coefficient (95\% CI 1.9 to 3.0; Holm-adjusted $P = 0.001$; lower error in one of the four regions). Within Alaska, with all labels, almost all of the gain of ML over the recalibrated model came from correcting bias between climate grid cells. The contribution is the evaluation design, its results including negative ones, and the label-count workflow built from them.
